% Regression: t2_pentagon3 — renders ^2F ^1F = ^1F ^2F ^2F as a physical string sliding through an MPO splitting…
% Target: RMP 2011.12127 fig3_pentagon3.pdf (pentagon-equation component).
% Formula:  ^2F ^1F = ^1F ^2F ^2F  —  a physical string (black, vertical,
% directed up) slides through an MPO splitting vertex (gray disc):
%   LHS: the string crosses the incoming red MPO wire WEST of the vertex;
%   RHS: the two split red MPO wires both cross the string EAST of it.
% Genuinely non-grid (free crossings + a trivalent vertex): tenkzfree.
% Honest gaps: \tnjoin has no wire-hue key, so the red (MPO) vs black
% (physical) distinction — the point of the figure — is carried only by
% geometry; the splitting vertex renders as a house bead, not a gray disc.
\documentclass[varwidth=19cm, border=6pt]{standalone}
\usepackage{amsmath, amssymb}
\usepackage{tenkz}
\begin{document}

\[
  \begin{tenkzfree}
    % the splitting vertex; incoming MPO wire from the west, two out east
    \tnput[dot, ports={west:virtual}]{V}{(0,0)}{}
    \tnjoin[dir]{-16mm,0mm}{V.west}
    \tnjoin[dir, route=curve, out=40, in=180]{V.north east}{15mm,8mm}
    \tnjoin[dir, route=curve, out=-40, in=180]{V.south east}{15mm,-8mm}
    % the physical string, west of the vertex, crossing the incoming wire
    \tnjoin[dir]{-9mm,-7mm}{-9mm,7mm}
  \end{tenkzfree}
  \;=\;
  \begin{tenkzfree}
    \tnput[dot, ports={west:virtual}]{W}{(0,0)}{}
    \tnjoin[dir]{-13mm,0mm}{W.west}
    \tnjoin[dir, route=curve, out=40, in=180]{W.north east}{19mm,8mm}
    \tnjoin[dir, route=curve, out=-40, in=180]{W.south east}{19mm,-8mm}
    % the string, now east of the vertex, crossing BOTH split wires
    \tnjoin[dir]{13mm,-13mm}{13mm,13mm}
  \end{tenkzfree}
\]
\[
  {}^{2}F\,{}^{1}F \;=\; {}^{1}F\,{}^{2}F\,{}^{2}F
\]

\end{document}
